% resume.tex
% vim:set ft=tex spell:

\documentclass[10pt,letterpaper]{article}
\usepackage[letterpaper,margin=0.75in]{geometry}
\usepackage[utf8]{inputenc}
\usepackage{mdwlist}
\usepackage[T1]{fontenc}
\usepackage{textcomp}
\usepackage{xcolor}
\usepackage{tgpagella}
\usepackage[colorlinks=true, linkcolor=black, urlcolor=black]{hyperref}
\usepackage{marvosym}
\usepackage{graphicx}
\usepackage{enumitem}

\pagestyle{empty}
\setlength{\tabcolsep}{0em}

% indentsection style, used for sections that aren't already in lists
% that need indentation to the level of all text in the document
\newenvironment{indentsection}[1]%
{\begin{list}{}%
	{\setlength{\leftmargin}{#1}}%
	\item[]%
}
{\end{list}}

% opposite of above; bump a section back toward the left margin
\newenvironment{unindentsection}[1]%
{\begin{list}{}%
	{\setlength{\leftmargin}{-0.5#1}}%
	\item[]%
}
{\end{list}}

% format two pieces of text, one left aligned and one right aligned
\newcommand{\headerrow}[2]
{\begin{tabular*}{\linewidth}{l@{\extracolsep{\fill}}r}
	#1 &
	#2 \\
\end{tabular*}}

% make "C++" look pretty when used in text by touching up the plus signs
\newcommand{\CPP}
{C\nolinebreak[4]\hspace{-.05em}\raisebox{.22ex}{\footnotesize\bf ++}}

% and the actual content starts here
\begin{document}

\begin{center}
{\LARGE \textbf{Joshua Kalpin}}
\vspace{0.1cm}

\Mundus \ \href{http://joshuakalp.in}{joshuakalp.in} \ \textbullet
\ \ \Letter \ \href{mailto:jkalpin@uwaterloo.ca}{jkalpin@uwaterloo.ca} \ \textbullet
\ \includegraphics[height=6pt]{Resources/Octocat} 
\href{https://github.com/kapin}{Kapin} \ \textbullet
 \includegraphics[height=10pt]{Resources/twitter}
\href{https://twitter.com/joshkalpin}{@joshkalpin}
\end{center}

\hrule
\vspace{-0.4em}
\subsection*{Technical Skills}

\begin{indentsection}{\parindent}
\hyphenpenalty=1000
\begin{description*}
	\item[Proficient Languages:]
	C, \CPP, Ruby, Java, Objective-C, C\#
	\item[Familiar Languages:]
	Python, Bash, JavaScript, \LaTeX, MATLAB, SQL
	\item[Frameworks and APIs:]
	Sinatra, iOS, JUnit, PostgreSQL, ASP.NET Web API
	\item[Tools:]
	Windows, OS X, Ubuntu, Git, Jenkins, GDB, Valgrind, Visual Studio
\end{description*}
\end{indentsection}

\hrule
\vspace{-0.4em}
\subsection*{Experience}

\begin{itemize}[leftmargin=*]
	\parskip=0.1em
	
	\item[]
	\headerrow
		{\textbf{SDE Intern - Azure Active Directory (AAD)}}
		{\textbf{Redmond, WA}}
	\\
	\headerrow
		{\emph{\href{https://www.microsoft.com}{Microsoft}}}
		{\emph{May 2014 - August 2014}}
	\begin{itemize*}
		\item Developed performance enhancements for the AAD Graph Api in C\#
		\item Re-architected the AAD Graph Api to use ASP.NET WebApi instead of WCF DS
		\item Researched the feasiblity of using an experimental tool to generate client SDKs
	\end{itemize*}	
	
	\item[]
	\headerrow
	{\textbf{Core Contributor}}
		{\textbf{Open Source}}
	\\
	\headerrow
		{\emph{\href{https://www.cocoapods.org}{CocoaPods}}}
		{\emph{October 2013 -- August 2014}}
	\begin{itemize*}
		\item CocoaPods is a package manager for iOS and OSX applications
		\item Developed enhancements for the linter to improve code coverage and quality
	\end{itemize*}		
	
	\item[]
	\headerrow
		{\textbf{Engineering Intern - iOS}}
		{\textbf{San Francisco, CA}}
	\\
	\headerrow
		{\emph{\href{https://www.lookout.com}{Lookout Inc.}}}
		{\emph{August 2013 -- December 2013}}
	\begin{itemize*}
		\item Designed a framework to provision iOS devices on specific carrier networks
		\item Developed a prototype of the iOS RESTful Contact Backup in a 24-hour hackathon
		\item Hacked the Android emulator to allow for new malware detection capabilities
	\end{itemize*}
	
	\item[]
	\headerrow
		{\textbf{Mobile Applications Developer - iOS}}
		{\textbf{Waterloo, ON}}
	\\
	\headerrow
		{\emph{\href{https://www.enflick.com}{Enflick Inc.}}}
		{\emph{December 2012 -- April 2013}}
	\begin{itemize*}
		\item Completed a technical rewrite and visual redesign of the TextNow iOS application
	\end{itemize*}

	\item[]
	\headerrow
		{\textbf{Ninja Software Developer}}
		{\textbf{Toronto, ON}}
	\\
	\headerrow
		{\emph{\href{https://www.work.com}{Salesforce.com - Work.com}}}
		{\emph{April -- August 2012}}
	\begin{itemize*}
		\item Developed both client- and server-side features in Java for the Work.com platform
	\end{itemize*}
\end{itemize}

\hrule
\vspace{-0.4em}
\subsection*{Projects}
	\begin{indentsection}{\parindent}
	\hyphenpenalty=1000
	\begin{description*}
	    \item[\href{https://github.com/SanaMobile/sana.protocol_builder}{Sana Protocol Builder}] A Web Application for developing medical procedures to use with the Sana Platform
		\item[\href{https://github.com/Kapin/paywall-smasher}{Paywall Smasher:}] A Google Chrome extension, written in JavaScript, that hides poorly designed paywalls
		\item[\href{https://github.com/Kapin/gimgur-server}{Gimgur Server:}] A small Sinatra app that creates Imgur albums when given a list of images
	\end{description*}
	\end{indentsection}

\hrule
\vspace{-0.4em}
\subsection*{Education}

\begin{itemize}[leftmargin=*]
	\parskip=0.1em

	\item[] 
	\headerrow
		{\textbf{University of Waterloo}}
		{\textbf{Waterloo, ON}}
	\\
	\headerrow
		{\emph{Candidate for a Bachelor of Software Engineering (BSE)}}
		{\emph{2011 -- 2016}}
		
	\item[] 
	\headerrow
		{\textbf{William Lyon Mackenzie Collegiate Institute}}
		{\textbf{North York, ON}}
	\\
	\headerrow
		{\emph{Ontario Secondary School Diploma}}
		{\emph{2007 -- 2011}}

\end{itemize}

\hrule
\vspace{-0.4em}
\subsection*{Activities and Interests}

\begin{itemize*}
\item Organizer of SE Hack Day, an event for creating and sharing cool projects
\item Vice President Education for the Waterloo Engineering Society
\item Budding audiophile and lover of a wide variety of music; especially progressive rock/metal and classical
\item Huge science fiction and fantasy nut with favourite authors including Brandon Sanderson and Isaac Asimov
\end{itemize*}



\end{document}
